% =============================================================================
% EBF WORKING PAPER TEMPLATE
% =============================================================================
%
% Format:     Working Paper
% Seiten:     15-30
% Aufwand:    2-4 Wochen
% Qualität:   Q2-Q3
% Ziel:       Forschung in Progress dokumentieren
% Leser:      Forscher, Team
%
% Nummerierung: WP-001, WP-002, ...
% Versionierung: v1, v2, v3, ...
%
% =============================================================================

\documentclass[12pt,a4paper]{article}
\usepackage[margin=2.5cm]{geometry}
\usepackage{amsmath,amssymb}
\usepackage{booktabs}
\usepackage{tcolorbox}
\usepackage{hyperref}
\usepackage{natbib}

% =============================================================================
% METADATEN (ANPASSEN!)
% =============================================================================
\newcommand{\WPNumber}{WP-001}
\newcommand{\WPVersion}{v1}
\newcommand{\WPDate}{Januar 2026}
\newcommand{\WPTitle}{[Titel des Working Papers]}
\newcommand{\WPAuthors}{[Autor 1], [Autor 2]}
\newcommand{\WPStatus}{Draft} % Draft / Under Review / Final

% =============================================================================
\begin{document}
% =============================================================================

% -----------------------------------------------------------------------------
% TITELSEITE
% -----------------------------------------------------------------------------
\begin{titlepage}
\begin{center}

\vspace*{2cm}

{\Large\textbf{BEATRIX Research Group}}\\[0.5cm]
{\large FehrAdvice \& Partners AG}\\[2cm]

{\huge\textbf{\WPTitle}}\\[1cm]

{\Large \WPAuthors}\\[2cm]

\begin{tcolorbox}[colback=blue!5!white, colframe=blue!75!black, width=8cm]
\centering
\textbf{Working Paper \WPNumber}\\
Version \WPVersion\\
\WPDate\\[0.5em]
Status: \textbf{\WPStatus}
\end{tcolorbox}

\vfill

{\small
\textbf{Abstract}\\[0.5em]
\begin{minipage}{0.8\textwidth}
[150-250 Wörter Abstract. Beschreibt Fragestellung, Methodik, 
Hauptergebnisse und Implikationen.]
\end{minipage}
}

\vspace{1cm}

{\footnotesize
Keywords: [keyword1], [keyword2], [keyword3]\\
JEL Codes: [D91], [C91], [...]
}

\end{center}
\end{titlepage}

% -----------------------------------------------------------------------------
% VERSION HISTORY
% -----------------------------------------------------------------------------
\begin{tcolorbox}[colback=gray!5!white, colframe=gray!50!black, 
    title=\textbf{Version History}]
\begin{tabular}{@{}llp{8cm}@{}}
\toprule
\textbf{Version} & \textbf{Datum} & \textbf{Änderungen} \\
\midrule
v1 & [Datum] & Initial draft \\
% v2 & [Datum] & [Beschreibung der Änderungen] \\
% v3 & [Datum] & [Beschreibung der Änderungen] \\
\bottomrule
\end{tabular}
\end{tcolorbox}

\tableofcontents
\newpage

% =============================================================================
\section{Introduction}
% =============================================================================

\subsection{Motivation}

[Warum ist diese Forschung wichtig? Was ist das Problem?]

\subsection{Research Question}

[Präzise Formulierung der Forschungsfrage]

\subsection{Contribution}

[Was trägt dieses Paper bei? Was ist neu?]

\subsection{Roadmap}

[Aufbau des Papers beschreiben]

% =============================================================================
\section{Related Literature}
% =============================================================================

[Einordnung in bestehende Literatur. Kurz und fokussiert.]

% =============================================================================
\section{Theoretical Framework}
% =============================================================================

\subsection{Setup}

[Grundlegende Annahmen und Definitionen]

\subsection{Model}

[Das theoretische Modell]

\subsection{Key Results}

[Wichtigste theoretische Ergebnisse]

% =============================================================================
\section{Empirical Strategy}
% =============================================================================

\subsection{Data}

[Datenbeschreibung]

\subsection{Methodology}

[Empirische Methodik]

% =============================================================================
\section{Results}
% =============================================================================

\subsection{Main Results}

[Hauptergebnisse]

\subsection{Robustness}

[Robustheitstests]

% =============================================================================
\section{Discussion}
% =============================================================================

\subsection{Interpretation}

[Interpretation der Ergebnisse]

\subsection{Limitations}

[Limitationen dieser Studie]

\subsection{Future Research}

[Offene Fragen und nächste Schritte]

% =============================================================================
\section{Conclusion}
% =============================================================================

[Zusammenfassung und Schlussfolgerungen]

% =============================================================================
% REFERENCES
% =============================================================================
\newpage
\bibliographystyle{apalike}
% \bibliography{references}

% Placeholder für Referenzen
\begin{thebibliography}{99}
\bibitem{placeholder} [Referenzen hier einfügen]
\end{thebibliography}

% =============================================================================
% APPENDIX
% =============================================================================
\appendix
\newpage
\section{Additional Results}

[Zusätzliche Tabellen, Abbildungen, Ableitungen]

\section{Data Appendix}

[Details zu Daten und Variablen]

\end{document}
